% =============================================================================
% Painel A1 - G1: LangGraph e Orquestracao Baseada em Grafos
% EN-045 - Agentes de IA
% Compilar com: pdflatex painel.tex  (rodar 2x)
% Recomendado: Overleaf (menu -> Compiler: pdfLaTeX)
% =============================================================================

\documentclass[25pt, a1paper, portrait, margin=20mm, innermargin=15mm,
               blockverticalspace=12mm, colspace=15mm, subcolspace=8mm]{tikzposter}

\usepackage[utf8]{inputenc}
\usepackage[T1]{fontenc}
\usepackage[brazil]{babel}
\usepackage{lmodern}
\usepackage{graphicx}
\usepackage{booktabs}
\usepackage{array}
\usepackage{xcolor}
\usepackage{tikz}
\usetikzlibrary{shapes.geometric, arrows.meta, positioning, fit, backgrounds}

% ---------- Tema --------------------------------------------------------------
\usetheme{Default}
\usecolorstyle{Britain}
\useblockstyle{Default}
\usetitlestyle{Default}

\definecolor{lgPrimary}{HTML}{1F3A93}
\definecolor{lgAccent}{HTML}{E67E22}
\definecolor{lgSoft}{HTML}{F4F6FB}

\colorlet{blocktitlebgcolor}{lgPrimary}
\colorlet{blocktitlefgcolor}{white}
\colorlet{blockbodybgcolor}{lgSoft}

% ---------- Metadados ---------------------------------------------------------
\title{LangGraph e Orquestração Baseada em Grafos}
\author{Cecilia Rufatto \quad Felipe Proença \quad Leonardo Moino \quad Leonardo Quartaroli \quad Lorena Scabello}
\institute{Agentes de IA \quad $|$ \quad Grupo G1}

\begin{document}

\maketitle

% =============================================================================
% LINHA 1 - Contexto + Conceitos-chave
% =============================================================================
\begin{columns}

\column{0.5}
\block{1. Contexto e motivação}{
Frameworks lineares como o LangChain tradicional executam etapas em ordem fixa. Funciona bem em fluxos curtos e determinísticos, mas agentes de verdade costumam precisar de coisas que esse modelo não entrega de forma nativa:

\begin{itemize}
  \item ciclos, pra refletir ou retentar uma etapa;
  \item estado persistente entre execuções;
  \item decisões condicionais que dependem do que aconteceu antes;
  \item pausa e retomada com intervenção humana.
\end{itemize}

\vspace{2pt}
O LangGraph cobre essa lacuna tratando o agente como um grafo: cada nó é uma função sobre o estado, as arestas decidem o caminho, e \textit{checkpointers} oficiais cuidam da persistência. É a base do LangSmith, do LangChain Academy e de agentes em produção em empresas como Klarna, Replit e Elastic.
}

\column{0.5}
\block{2. Conceitos-chave}{
\begin{tabular}{@{} >{\bfseries}p{0.30\linewidth} p{0.65\linewidth} @{}}
\toprule
StateGraph & Grafo cujo estado é um \texttt{TypedDict} compartilhado entre nós. \\
Node & Função Python que recebe e devolve estado. \\
Edge & Conexão direta entre dois nós (\texttt{add\_edge}). \\
Conditional Edge & Roteamento dinâmico via função \texttt{router(state)}. \\
Reducer & Estratégia de \emph{merge} de campos do estado (ex.\ \texttt{add\_messages}). \\
Checkpointer & Persistência do estado entre execuções (SQLite, Postgres, Redis). \\
HITL & \textit{Human-in-the-loop}: pausa com \texttt{interrupt} e retoma via \texttt{Command}. \\
\bottomrule
\end{tabular}
}

\end{columns}

% =============================================================================
% LINHA 2 - Estado da arte + Arquitetura
% =============================================================================
\begin{columns}

\column{0.5}
\block{3. Estado da arte}{
\textbf{Linha do tempo}\\[2pt]
\begin{tabular}{@{} p{0.18\linewidth} p{0.78\linewidth} @{}}
\toprule
\textbf{2022} & LangChain --- \textit{chains} lineares para composição de LLMs. \\
\textbf{Jan 2024} & LangGraph 0.0.x --- primeira versão pública. \\
\textbf{Mai 2024} & LangGraph Studio + Cloud --- observabilidade e deploy. \\
\textbf{2025} & LangGraph 1.x estável --- \emph{checkpointers}, durable execution. \\
\bottomrule
\end{tabular}

\vspace{6pt}
\textbf{Comparativo com frameworks concorrentes}\\[2pt]
\begin{tabular}{@{} p{0.22\linewidth} p{0.20\linewidth} p{0.20\linewidth} p{0.28\linewidth} @{}}
\toprule
\textbf{Framework} & \textbf{Estado} & \textbf{Ciclos} & \textbf{Persistência} \\
\midrule
LangGraph    & Tipado    & Nativo  & Checkpointers oficiais \\
CrewAI       & Implícito & Limitado & Memory module \\
AutoGen      & Mensagens & Por chat & Externa \\
OpenAI Swarm & Handoffs  & Não     & Não (experimental) \\
\bottomrule
\end{tabular}
}

\column{0.5}
\block{4. Arquitetura do protótipo}{
\centering
\includegraphics[width=0.95\linewidth]{metrics/graph.png}\\[6pt]
\raggedright
Cada mensagem que entra é processada pelo nó \texttt{detect\_intent}, que rotula como \textit{pedido}, \textit{suporte} ou \textit{geral}. A aresta condicional \texttt{route\_intent} despacha pro handler correspondente, que responde e grava o histórico em JSON pelo \texttt{memory.py} antes de chegar ao \texttt{END}.\\[4pt]

O estado é um \texttt{TypedDict} (\texttt{AgentState}) com os campos \texttt{message, intent, confidence, response, history, metadata, timestamp} e \texttt{session\_id}.
}

\end{columns}

% =============================================================================
% LINHA 3 - Experimento + Discussao
% =============================================================================
\begin{columns}

\column{0.5}
\block{5. Experimento e resultados}{
Montamos um dataset de 40 mensagens em português rotuladas em três intenções (\texttt{pedido}=15, \texttt{suporte}=13, \texttt{geral}=12). Cada mensagem passa pelo grafo compilado; medimos latência com \texttt{time.perf\_counter} e calculamos as métricas com \texttt{scikit-learn}.\\[4pt]

\begin{minipage}[t]{0.48\linewidth}
\centering
\textbf{Métricas globais}\\[2pt]
\begin{tabular}{@{} l r @{}}
\toprule
Accuracy        & \textbf{85{,}00\%} \\
Macro F1        & 0{,}8507 \\
Latência média  & 33{,}52 ms \\
Latência P95    & $\sim$56 ms \\
\bottomrule
\end{tabular}

\vspace{6pt}
\textbf{Por classe}\\[2pt]
\begin{tabular}{@{} l r r r @{}}
\toprule
Classe   & P     & R     & F1 \\
\midrule
pedido   & 0{,}92 & 0{,}80 & 0{,}86 \\
suporte  & 0{,}92 & 0{,}85 & 0{,}88 \\
geral    & 0{,}73 & 0{,}92 & 0{,}81 \\
\bottomrule
\end{tabular}
\end{minipage}\hfill
\begin{minipage}[t]{0.48\linewidth}
\centering
\includegraphics[width=\linewidth]{metrics/confusion_matrix.png}\\
\small Matriz de confusão (n=40)
\end{minipage}
}

\column{0.5}
\block{6. Discussão e limitações}{
\textbf{O que funcionou.} A combinação de StateGraph com aresta condicional deixou o roteamento declarativo e fácil de testar --- conseguimos escrever evals direto sobre o grafo compilado. Persistência em JSON foi suficiente para o protótipo, e a latência abaixo de 60 ms no P95 mantém o agente confortável dentro de um chat síncrono.\\[4pt]

\textbf{Limitações.} Classificar por keywords não captura semântica: a classe \texttt{geral} acaba absorvendo qualquer frase que não casa explicitamente com as outras --- daí o \emph{precision} de 0{,}73. Mensagens multi-intenção, como ``meu produto chegou quebrado e quero suporte'', são forçadas a uma única rota. Também não há suporte a contexto profundo entre turnos.\\[4pt]

\textbf{Próximos passos.} Trocar o \texttt{detect\_intent} por um classificador LLM \emph{few-shot} com saída estruturada via Pydantic; adicionar um nó de Human-in-the-Loop para casos em que o \texttt{confidence} fique baixo; e migrar a persistência JSON para um \texttt{SqliteSaver}, ganhando \emph{durable execution}.
}

\end{columns}

% =============================================================================
% RODAPE - Referencias + QR Code
% =============================================================================
\block[bodyoffsety=10mm, titleoffsety=10mm]{Referências e repositório}{
\begin{minipage}[c]{0.78\linewidth}
\footnotesize
\begin{enumerate}
  \item LANGCHAIN AI. \textbf{LangGraph --- Overview and Graph API}. Documentação oficial. Disponível em: \texttt{langchain-ai.github.io/langgraph/}.
  \item LANGCHAIN AI. \textbf{Persistence and Checkpointers}. \texttt{langchain-ai.github.io/langgraph/concepts/persistence/}.
  \item LANGCHAIN AI. \textbf{Human-in-the-loop}. \texttt{langchain-ai.github.io/langgraph/concepts/human\_in\_the\_loop/}.
  \item WU, Q. \emph{et al.} \textbf{AutoGen: Enabling Next-Gen LLM Applications via Multi-Agent Conversation}. arXiv:2308.08155, 2023.
  \item YAO, S. \emph{et al.} \textbf{ReAct: Synergizing Reasoning and Acting in Language Models}. ICLR, 2023.
  \item OPENAI. \textbf{Swarm}: educational framework for multi-agent orchestration. GitHub, 2024.
  \item LANGCHAIN AI. \textbf{LangGraph Academy} --- curso oficial. \texttt{academy.langchain.com/courses/intro-to-langgraph}.
\end{enumerate}
\end{minipage}\hfill
\begin{minipage}[c]{0.20\linewidth}
\centering

\includegraphics[width=0.85\linewidth]{metrics/qrcode_repo.png}\\
\footnotesize \textbf{Repositório do grupo}\\
\texttt{github.com/Ceciliarufatto/langgraph}
\end{minipage}
}

\end{document}
